% !Mode:: "TeX:UTF-8"
% !TEX program  = xelatex

\section{相关工作}
\label{sec:related}

\subsection{神经组合优化与方法选择}

神经组合优化使用学习模型直接构造组合优化问题的求解策略。
早期工作采用强化学习训练序列解码器，直接生成组合优化问题的解\cite{bello2016nco}，
随后基于注意力机制的方法进一步提高了推理效率和求解质量\cite{kool2019attention}。
不过，不同方法在不同实例分布上往往各有优势。

算法选择问题可以追溯到 Rice 的经典框架：
算法性能取决于实例特征与算法特性的匹配关系\cite{rice1976algorithmselection}。
对于 TSP 与 CVRP，不同神经方法受到网络结构、训练分布和推理机制的影响，在实例层面具有互补性。

Gao et al.\cite{gao2025nss}将这一思路引入神经组合优化领域，
提出 NSS 框架。该工作通过逐实例对比确认了方法之间的互补性，并把问题拆分为特征提取、选择模型和选择策略三个部分，
说明实例级方法选择确实是一个成立的问题。
与NSS 相比，本文在两个方面有所不同：
一，NSS 采用全信息监督学习，而本文采用在线的上下文老虎机更新，更贴近逐实例部分反馈的实际场景；
二，NSS 对TSP 和 CVRP 分别训练独立模型，而本文用单一模型统一处理两类问题，通过共享动作空间与按问题拆分的浅层头实现跨问题的经验共享。

\subsection{上下文老虎机与 Neural-LinUCB}

上下文老虎机在每一轮观测上下文后从动作集合中选择一个动作并获得即时反馈，不涉及多步状态转移。
LinUCB 利用线性模型估计期望奖励，通过上置信界（Upper Confidence Bound, UCB）
平衡探索与利用\cite{auer2002ucb1,li2010linucb}。

Neural-LinUCB\cite{xu2022neurallinucb}将深度表征学习与浅层线性探索分开处理：
用神经网络学习实例表示，只在最后的低维表示空间中维护 LinUCB 统计量，并周期性更新隐藏层。
这一思路与本文任务契合，因为算子选择本身就是单步决策，关键在于学到有区分度的实例表示，并在低维空间内保持稳定探索。

综上，现有工作已初步验证了实例级方法选择的可行性，
也提供了上下文老虎机与深度表征结合的理论方法。
但在神经组合优化场景下，如何用单一模型同时服务多类问题、如何在在线部分反馈设定下学习选择策略，仍有待进一步探索。
